% Wersja źródłowa: audyt II, jednostka 10
\setcounter{section}{9}
\section{Przestrzenie beczkowe i bornologiczne. Twierdzenie Banacha--Steinhausa}\label{sec:beczkowe-bornologiczne}

\subsection{Od ograniczoności punktowej do równociągłości}

W rozdziale o twierdzeniu Banacha--Alaoglu pokazaliśmy, że rodzina
$H\subseteq E'$ jest równociągła dokładnie wtedy, gdy mieści się
w zbiorze polarnym pewnego otoczenia zera. Równociągłość implikuje
ograniczoność punktową, lecz przykład z funkcjonałami $f_n(x)=nx_n$ na $c_{00}$ pokazuje, że implikacji nie można odwrócić w każdej
przestrzeni liniowo-topologicznej. Zbadamy teraz, jakie dodatkowe
własności przestrzeni dziedziny gwarantują tę implikację.

W całym rozdziale $\K\in\{\R,\C\}$. Przestrzenie $E$ i $F$
są lokalnie wypukłymi przestrzeniami liniowo-topologicznymi Hausdorffa,
chyba że wyraźnie zaznaczono inaczej. Symbol $E'$ oznacza przestrzeń
dualną topologiczną, a $\mathcal L(E,F)$ zbiór ciągłych operatorów
liniowych z $E$ do $F$. Przez zbiór ograniczony rozumiemy zbiór
ograniczony w sensie von Neumanna. Zbiór $C$ jest \emph{absolutnie
wypukły}, jeżeli jest wypukły i zrównoważony.

\begin{definicja}[Beczka i przestrzeń beczkowa]\label{def:beczka}
Niech $E$ będzie lokalnie wypukłą TVS. Zbiór $B\subseteq E$
nazywamy \emph{beczką}, jeżeli jest domknięty, absolutnie wypukły
i pochłaniający, to znaczy dla każdego $x\in E$ istnieje $t>0$
takie, że $x\in tB$. Przestrzeń $E$ jest \emph{beczkowa}, jeśli
każda beczka w $E$ jest otoczeniem zera.
\end{definicja}

Warunek pochłaniania dotyczy pojedynczych punktów, a nie całych
zbiorów ograniczonych. To rozróżnienie będzie istotne przy
wprowadzaniu przestrzeni bornologicznych.

\begin{lemat}[Zbiór polarny i ograniczoność punktowa]\label{lem:polar-barrel}
Niech $H\subseteq E'$ i niech
$H^\circ=\{x\in E:|f(x)|\le1\ \text{dla każdego }f\in H\}$.
Wówczas $H^\circ$ jest domknięty i absolutnie wypukły.
Ponadto $H^\circ$ jest pochłaniający wtedy i tylko wtedy, gdy
$H$ jest punktowo ograniczony, czyli
$\sup_{f\in H}|f(x)|<\infty$ dla każdego $x\in E$.
\end{lemat}
\begin{proof}
Każdy zbiór $\{x:|f(x)|\le1\}$ jest domknięty i absolutnie
wypukły, a przecięcie dowolnej rodziny takich zbiorów
zachowuje obie własności. Jeśli $H$ jest punktowo ograniczony,
to dla ustalonego $x$ liczba $M_x=\sup_{f\in H}|f(x)|$
jest skończona; przyjmujemy $M_x=0$, gdy $H=\varnothing$.
Dla $t>\max\{1,M_x\}$ mamy $x/t\in H^\circ$, czyli
$x\in tH^\circ$. Odwrotnie, jeśli $x\in tH^\circ$, to
$|f(x)|\le t$ dla wszystkich $f\in H$. W obu kierunkach
dopuszczamy również $x=0$ i rodzinę pustą.
\end{proof}

\begin{twierdzenie}[Charakteryzacja beczkowości]\label{thm:barrelled-char}
Niech $E$ będzie lokalnie wypukłą TVS Hausdorffa.
Następujące warunki są równoważne:
\begin{enumerate}[label=\textup{(\roman*)}]
\item $E$ jest beczkowa;
\item każda punktowo ograniczona rodzina $H\subseteq E'$
      jest równociągła;
\item każda rodzina $H\subseteq E'$ ograniczona w
      $\sigma(E',E)$ jest równociągła.
\end{enumerate}
\end{twierdzenie}
\begin{proof}
Warunki (ii) i (iii) są równoważne, gdyż topologia
$\sigma(E',E)$ jest topologią punktowej zbieżności,
a ograniczoność zbioru w produkcie jest równoważna
ograniczoności każdej jego projekcji na $\K$.

Załóżmy (i) i niech $H$ będzie punktowo ograniczona.
Z lematu~\ref{lem:polar-barrel} wynika, że $H^\circ$
jest beczką, zatem jest otoczeniem zera. Dla
$u\in H^\circ$ zachodzi $|f(u)|\le1$ dla wszystkich
$f\in H$, czyli $H$ jest równociągła.

Załóżmy (ii) i niech $B\subseteq E$ będzie beczką.
Zbiór polarny $B^\circ\subseteq E'$ jest punktowo ograniczony:
dla każdego $x$ istnieje $t>0$ z $x\in tB$, zatem
$|f(x)|\le t$ dla wszystkich $f\in B^\circ$.
Z (ii) otrzymujemy otoczenie $U$ takie, że
$B^\circ\subseteq U^\circ$. Wskutek odwracania inkluzji
przez operację przejścia do zbioru polarnego i twierdzenia
bipolarnego dla lokalnie wypukłej przestrzeni Hausdorffa
zachodzi
$U\subseteq U^{\circ\circ}\subseteq B^{\circ\circ}=B$.
Zatem $B$ jest otoczeniem zera.
\end{proof}

\begin{uwaga}[Dlaczego konieczna jest lokalna wypukłość?]
W dowodzie odwrotnej implikacji użyliśmy równości
$B^{\circ\circ}=B$ dla domkniętego zbioru absolutnie
wypukłego. Opiera się ona na twierdzeniu Hahna--Banacha
i może zawieść w ogólnej TVS. Charakteryzację
\ref{thm:barrelled-char} formułujemy więc w klasie
przestrzeni lokalnie wypukłych Hausdorffa.
\end{uwaga}

\subsection{Własność Baire'a i przestrzenie Frécheta}

\begin{definicja}[Przestrzeń Baire'a]
Przestrzeń topologiczną $X$ nazywamy \emph{przestrzenią Baire'a},
jeżeli każde przeliczalne przecięcie otwartych zbiorów gęstych
w $X$ jest gęste. Równoważnie: przeliczalna suma zbiorów
domkniętych o pustym wnętrzu ma puste wnętrze.
\end{definicja}

\begin{twierdzenie}[Baire: zupełność metryczna wystarcza]
\label{thm:baire-complete}
Każda niepusta zupełna przestrzeń metryczna $(X,d)$
jest przestrzenią Baire'a.
\end{twierdzenie}
\begin{proof}
Niech $G_1,G_2,\ldots$ będą otwartymi podzbiorami
gęstymi w $X$. Pokażemy, że ich przecięcie przecina
każdy niepusty zbiór otwarty $O\subseteq X$.
Wybieramy indukcyjnie punkty $x_n\in X$ oraz
liczby $0<r_n<2^{-n}$ tak, by
\[
 \overline{B}(x_1,r_1)\subseteq O\cap G_1,\qquad
 \overline{B}(x_{n+1},r_{n+1})\subseteq
 B(x_n,r_n)\cap G_{n+1}.
\]
W pierwszym kroku $O\cap G_1$ jest niepusty i otwarty,
więc zawiera domkniętą kulę o dodatnim promieniu.
Analogicznie $B(x_n,r_n)\cap G_{n+1}$ jest niepusty
i otwarty, ponieważ $G_{n+1}$ jest gęsty. Możemy
zatem kontynuować konstrukcję.

Dla $m>n$ mamy $x_m\in B(x_n,r_n)$, więc
$d(x_m,x_n)<r_n<2^{-n}$. Ciąg $(x_n)$ jest ciągiem
Cauchy'ego. Z zupełności otrzymujemy punkt $x\in X$
z $x_n\to x$. Dla ustalonego $n$ wszystkie punkty
$x_m$ o $m>n$ leżą w $\overline B(x_n,r_n)$,
a domkniętość tej kuli daje
$x\in\overline B(x_n,r_n)\subseteq G_n$.
Ponadto $x\in\overline B(x_1,r_1)\subseteq O$.
Wobec dowolności $O$ przecięcie $\bigcap_nG_n$
jest gęste w $X$.
\end{proof}

\begin{twierdzenie}[Przestrzenie Baire'a są beczkowe]\label{thm:baire-barrelled}
Niech $E$ będzie lokalnie wypukłą TVS Hausdorffa, której
przestrzeń topologiczna jest przestrzenią Baire'a.
Wówczas $E$ jest beczkowa.
\end{twierdzenie}
\begin{proof}
Niech $B$ będzie beczką w $E$. Ponieważ jest pochłaniająca
i zrównoważona, $E=\bigcup_{n=1}^\infty nB$. Każdy zbiór
$nB$ jest domknięty, bo mnożenie przez niezerowy skalar
jest homeomorfizmem. Z własności Baire'a otrzymujemy
$\operatorname{int}(n_0B)\ne\varnothing$ dla pewnego
$n_0\ge1$, a zatem $\operatorname{int}B\ne\varnothing$.

Wybierzmy $y\in\operatorname{int}B$. Zrównoważenie daje
$-y\in\operatorname{int}B$. Istnieje zrównoważone
otwarte otoczenie $V$ zera takie, że
$y+V\subseteq B$ i $-y+V\subseteq B$. Dla każdego
$v\in V$ wypukłość $B$ daje
$v=\frac12(y+v)+\frac12(-y+v)\in B$.
Zatem $V\subseteq B$ i $B$ jest otoczeniem zera.
\end{proof}

\begin{wniosek}[Przestrzenie Frécheta i Banacha]\label{cor:frechet-barrelled}
Każda przestrzeń Frécheta i każda przestrzeń Banacha
jest beczkowa. Również każda lokalnie wypukła TVS
Hausdorffa, której przestrzeń topologiczna jest Baire'a,
jest beczkowa.
\end{wniosek}
\begin{proof}
Z Twierdzenia~\ref{thm:baire-complete} wynika, że
przestrzeń Frécheta jest Baire'a.
Przestrzeń Banacha jest szczególnym przypadkiem przestrzeni
Frécheta. Pozostała część jest
Twierdzeniem~\ref{thm:baire-barrelled}.
\end{proof}

\begin{uwaga}[Dwa twierdzenia o nazwie Baire'a]
Należy odróżnić własność Baire'a przestrzeni topologicznej
od twierdzenia Baire'a, które pozwala ją uzyskać z
zupełności metrycznej. Dowód twierdzenia Baire'a podaliśmy powyżej;
w następnym rozdziale zastosujemy go do twierdzeń
o odwzorowaniu otwartym i wykresie domkniętym.
\end{uwaga}

\subsection{Twierdzenie Banacha--Steinhausa w przestrzeniach lokalnie wypukłych}

\begin{definicja}[Ograniczoność punktowa i równociągłość operatorów]
Niech $\mathcal T\subseteq\mathcal L(E,F)$.
Rodzina $\mathcal T$ jest \emph{punktowo ograniczona}, gdy
$\{Tx:T\in\mathcal T\}$ jest ograniczony w $F$ dla każdego
$x\in E$. Jest \emph{równociągła}, gdy dla każdego otoczenia
$V$ zera w $F$ istnieje otoczenie $U$ zera w $E$
takie, że $T(U)\subseteq V$ dla każdego $T\in\mathcal T$.
\end{definicja}

\begin{twierdzenie}[Banach--Steinhaus dla przestrzeni beczkowych]
\label{thm:banach-steinhaus-lcs}
Niech $E$ będzie beczkową lokalnie wypukłą TVS Hausdorffa,
$F$ lokalnie wypukłą TVS Hausdorffa, a
$\mathcal T\subseteq\mathcal L(E,F)$ rodziną operatorów.
Jeżeli $\mathcal T$ jest punktowo ograniczona, to jest
równociągła.
\end{twierdzenie}
\begin{proof}
Ustalmy otoczenie $V$ zera w $F$. Z lokalnej wypukłości
możemy wybrać ciągłe seminormy $q_1,\ldots,q_m$ na $F$
i $\varepsilon>0$ takie, aby
$W=\{y:\max_jq_j(y)\le\varepsilon\}\subseteq V$.
Zbiór $W$ jest domknięty, absolutnie wypukły i stanowi
otoczenie zera w $F$. Określmy
\[
 B=\bigcap_{T\in\mathcal T}T^{-1}(W).
\]
Jest to zbiór domknięty i absolutnie wypukły w $E$.
Dla ustalonego $x\in E$ zbiór $\{Tx:T\in\mathcal T\}$
jest ograniczony w $F$. Istnieje więc $t_x>0$ takie,
że $Tx\in t_x W$ dla każdego $T\in\mathcal T$,
a zatem $x\in t_xB$. Zbiór $B$ jest beczką.
Beczkowość $E$ daje otoczenie zera $U\subseteq B$,
więc $T(U)\subseteq W\subseteq V$ dla wszystkich
$T\in\mathcal T$.
\end{proof}

\begin{wniosek}[Klasyczna zasada jednostajnej ograniczoności]
\label{cor:uniform-boundedness}
Niech $E$ będzie przestrzenią Banacha, a $F$ przestrzenią
unormowaną nad $\K$. Jeżeli rodzina
$\mathcal T\subseteq\mathcal L(E,F)$ spełnia
$\sup_{T\in\mathcal T}\|Tx\|_F<\infty$ dla każdego
$x\in E$, to $\sup_{T\in\mathcal T}\|T\|<\infty$.
Nie zakładamy zupełności $F$.
\end{wniosek}
\begin{proof}
Przestrzeń $E$ jest beczkowa. Stosujemy
Twierdzenie~\ref{thm:banach-steinhaus-lcs} do kuli
$V=\{y:\|y\|_F<1\}$. Istnieje $\delta>0$ takie, że
$\|Tx\|_F<1$ dla wszystkich $T\in\mathcal T$
i $\|x\|_E<\delta$. Przez skalowanie otrzymujemy
$\|T\|\le 2/\delta$ dla każdego $T\in\mathcal T$.
\end{proof}

\begin{proof}[Drugi dowód wniosku: bezpośrednio z twierdzenia Baire'a]
Dla $n\in\N$ niech
$A_n=\{x\in E:\sup_{T\in\mathcal T}\|Tx\|_F\le n\}$.
Są to zbiory domknięte, ponieważ każdy $T$ jest ciągły.
Ograniczoność punktowa oznacza
$E=\bigcup_{n=1}^\infty A_n$. Na mocy twierdzenia Baire'a
pewien $A_N$ ma niepuste wnętrze. Istnieją więc
$x_0\in E$ i $r>0$ takie, że
$x_0+\{h:\|h\|_E<r\}\subseteq A_N$.
Dla $\|h\|_E<r$ i każdego $T\in\mathcal T$ mamy
$\|Th\|_F\le\|T(x_0+h)\|_F+\|Tx_0\|_F\le2N$.
Skalowanie daje $\|T\|\le4N/r$, niezależnie od $T$.
W dowodzie nie wystąpiła zupełność przestrzeni wartości $F$.
\end{proof}

\begin{przyklad}[Dlaczego zupełność dziedziny jest istotna?]
\label{ex:c00-not-barrelled}
Niech $E=c_{00}$ z normą $\|x\|_\infty=\sup_n|x_n|$.
Dla $n\ge1$ określmy $f_n(x)=nx_n$. Każdy $f_n\in E'$,
ponieważ $|f_n(x)|\le n\|x\|_\infty$.
Dla ustalonego $x\in c_{00}$ ciąg $(f_n(x))$ jest
ostatecznie zerowy, więc rodzina $H=\{f_n:n\ge1\}$ jest
punktowo ograniczona. Jednak $\|f_n\|=n\to\infty$,
zatem rodzina nie jest równociągła.

Wprost widać też beczkę niebędącą otoczeniem zera:
\[
 B=\{x\in c_{00}: n|x_n|\le1\ \text{dla każdego }n\ge1\}.
\]
Jest ona domknięta, absolutnie wypukła i pochłaniająca,
gdyż każdy $x$ ma skończony nośnik. Nie jest otoczeniem
zera: dla $x^{(n)}=n^{-1/2}e_n$ zachodzi
$\|x^{(n)}\|_\infty\to0$, podczas gdy
$n|x_n^{(n)}|=\sqrt n>1$ dla $n>1$. Tym samym
$E$ nie jest beczkowa.
\end{przyklad}

\begin{przyklad}[Topologia słaba na przestrzeni Banacha]
\label{ex:weak-not-barrelled}
Niech $E$ będzie nieskończenie wymiarową przestrzenią Banacha
i wyposażmy ją w topologię $\sigma(E,E')$. Domknięta kula
normowa $B=\{x:\|x\|\le1\}$ jest słabo domknięta
(z twierdzenia Hahna--Banacha), absolutnie wypukła i pochłaniająca.
Nie jest jednak słabym otoczeniem zera, gdyż każde słabe
otoczenie zawiera niezerową prostą, która nie mieści się
w $B$. Zatem $(E,\sigma(E,E'))$ nie jest beczkowa,
mimo że $E$ jest zupełna w swojej oryginalnej normie.
O beczkowości decyduje \emph{wybrana topologia}, a nie tylko
struktura algebraiczna przestrzeni.
\end{przyklad}

\subsection{Trzy konsekwencje zasady jednostajnej ograniczoności}

\begin{twierdzenie}[Punktowa zbieżność operatorów]\label{thm:pointwise-limit-cont}
Niech $E$ będzie beczkową lokalnie wypukłą TVS Hausdorffa,
a $F$ lokalnie wypukłą TVS Hausdorffa. Załóżmy, że
$(T_i)_{i\in I}$ jest siecią w $\mathcal L(E,F)$,
rodzina $\{T_i:i\in I\}$ jest punktowo ograniczona,
a dla każdego $x\in E$ istnieje granica
$T(x)=\lim_iT_i(x)$ w $F$. Wówczas $T:E\to F$
jest ciągłym operatorem liniowym, a rodzina
$\{T_i:i\in I\}$ jest równociągła.
\end{twierdzenie}
\begin{proof}
Liniowość $T$ wynika z ciągłości dodawania i mnożenia
przez skalar w $F$. Z ograniczoności punktowej i
Twierdzenia~\ref{thm:banach-steinhaus-lcs} otrzymujemy
równociągłość rodziny. Dla danego otoczenia $V$ zera
w $F$ wybieramy domknięte otoczenie $W\subseteq V$.
Istnieje otoczenie $U$ zera w $E$ takie, że
$T_i(U)\subseteq W$ dla wszystkich $i$.
Przechodząc do granicy przy ustalonym $x\in U$, otrzymujemy
$T(x)\in W\subseteq V$. Wobec dowolności $V$ operator
$T$ jest ciągły w zerze, a z liniowości także na całym $E$.
\end{proof}

\begin{wniosek}[Granica punktowa ciągu operatorów]
\label{cor:sequence-operators}
Niech $E$ i $F$ spełniają założenia
Twierdzenia~\ref{thm:pointwise-limit-cont}, a
$T_n\in\mathcal L(E,F)$ dla $n\ge1$. Jeśli
$T_n(x)\to T(x)$ dla każdego $x\in E$, to $T$ jest
ciągłym operatorem liniowym, a rodzina $\{T_n:n\ge1\}$
jest równociągła.
\end{wniosek}
\begin{proof}
Dla każdego $x\in E$ zbiór wartości zbieżnego ciągu
$(T_n(x))$ jest ograniczony w $F$ (jego ogon zbiega do zera
po odjęciu granicy, a skończony początek jest ograniczony).
Stosujemy Twierdzenie~\ref{thm:pointwise-limit-cont}.
\end{proof}

\begin{uwaga}[Dlaczego nie wystarcza sama zbieżność sieci?]
Dla sieci zbieżnej zbiór wszystkich jej wartości nie musi
być ograniczony. Niech $I$ oznacza zbiór skończonych
podzbiorów $\N$, uporządkowany inkluzją, i niech
$a_F=|F|$, gdy $1\notin F$, oraz $a_F=0$, gdy $1\in F$.
Wtedy $a_F\to0$ jako sieć w $\R$, ponieważ na całym
ogonie zaczynającym się od $\{1\}$ mamy $a_F=0$.
Jednocześnie zbiór wszystkich $a_F$ jest nieograniczony.
Przyjmując $E=F=\R$ i $T_F(x)=a_Fx$, otrzymujemy sieć
ciągłych operatorów liniowych zbieżną punktowo do zera,
która nie jest równociągła. Dlatego w twierdzeniu dla
ogólnych sieci jawnie zakładamy ograniczoność punktową.
\end{uwaga}

\begin{twierdzenie}[Jednostajność zbieżności na zbiorach zwartych]
\label{thm:equicont-compact-conv}
Niech $E$ i $F$ będą TVS Hausdorffa,
$\mathcal T\subseteq\mathcal L(E,F)$ rodziną równociągłą,
a $(T_i)_{i\in I}$ siecią z $\mathcal T$ zbieżną punktowo
do operatora $T\in\mathcal L(E,F)$.
Wówczas $T_i\to T$ jednostajnie na każdym zwartym
zbiorze $K\subseteq E$: dla każdego otoczenia $V$ zera
w $F$ istnieje $i_0\in I$ takie, że
$(T_i-T)(K)\subseteq V$ dla $i\ge i_0$.
\end{twierdzenie}
\begin{proof}
Ustalmy $K$ i otoczenie $V$. Wybierzmy zrównoważone
otoczenie $W$ z $W+W+W\subseteq V$. Równociągłość
rodziny $\mathcal T$ i ciągłość $T$ dają wspólne
otoczenie zera $U$ w $E$, które możemy wybrać
zrównoważone, tak że $T_i(U)\subseteq W$ dla wszystkich
$i$ oraz $T(U)\subseteq W$.
Wybieramy $U$ dodatkowo otwarte. Zwartość $K$
pozwala pokryć go skończenie wieloma zbiorami
$x_1+U,\ldots,x_m+U$ o środkach $x_j\in K$.
Punktowa zbieżność zapewnia indeks $i_0$ z
$(T_i-T)x_j\in W$ dla wszystkich $j$ i $i\ge i_0$.
Jeśli $x\in K\cap(x_j+U)$, to
$(T_i-T)x=(T_i-T)x_j+T_i(x-x_j)-T(x-x_j)\in W+W+W
\subseteq V$, co kończy dowód.
\end{proof}

\begin{stwierdzenie}[Ograniczoność w topologiach operatorowych]
\label{prop:sot-bot-bounded}
Niech $E$ będzie beczkową lokalnie wypukłą TVS Hausdorffa,
a $F$ lokalnie wypukłą TVS Hausdorffa. Na $\mathcal L(E,F)$
rozważmy topologię $\tau_s$ zbieżności punktowej i
$\tau_b$ zbieżności jednostajnej na zbiorach ograniczonych.
Wtedy rodzina $\mathcal T\subseteq\mathcal L(E,F)$
jest ograniczona w $\tau_s$ wtedy i tylko wtedy,
gdy jest ograniczona w $\tau_b$.
\end{stwierdzenie}
\begin{proof}
Ponieważ singletony są ograniczone, $\tau_b$ jest
silniejsza od $\tau_s$, więc ograniczoność w $\tau_b$
implikuje ograniczoność w $\tau_s$.

Załóżmy odwrotnie, że $\mathcal T$ jest ograniczona
w $\tau_s$. Zbieżność punktowa oznacza, że dla każdego
$x\in E$ zbiór $\{Tx:T\in\mathcal T\}$ jest ograniczony
w $F$. Z twierdzenia Banacha--Steinhausa rodzina jest
równociągła. Dla ograniczonego $B\subseteq E$ i otoczenia
$V$ zera w $F$ istnieje otoczenie $U$ zera w $E$
takie, że $T(U)\subseteq V$ dla wszystkich $T$.
Ponieważ $B$ jest ograniczony, istnieje $t>0$
z $B\subseteq tU$. Stąd $T(B)\subseteq tV$ dla wszystkich
$T\in\mathcal T$. Oznacza to ograniczoność w $\tau_b$.
\end{proof}

\subsection{Przestrzenie bornologiczne}

\begin{definicja}[Zbiór pochłaniający zbiory ograniczone]
\label{def:bornivorous}
Niech $E$ będzie TVS. Zbiór $C\subseteq E$ nazywamy
\emph{pochłaniającym zbiory ograniczone}, jeśli dla każdego
ograniczonego $B\subseteq E$ istnieje $t>0$ takie, że
$B\subseteq tC$. W szczególności $C$ jest pochłaniający,
ponieważ wszystkie singletony są ograniczone.
\end{definicja}

\begin{definicja}[Przestrzeń bornologiczna]\label{def:bornological}
Lokalnie wypukłą TVS Hausdorffa $E$ nazywamy
\emph{bornologiczną}, jeżeli każdy absolutnie wypukły
zbiór pochłaniający zbiory ograniczone w $E$ jest
otoczeniem zera. Nie wymagamy jego domkniętości.
\end{definicja}

\begin{lemat}[Funkcjonał Minkowskiego i pochłanianie ograniczonych]
\label{lem:minkowski-bornivorous}
Niech $C\subseteq E$ będzie absolutnie wypukły i
pochłaniający. Jego funkcjonał Minkowskiego
$p_C(x)=\inf\{t>0:x\in tC\}$ jest skończoną seminormą.
Zbiór $C$ pochłania wszystkie zbiory ograniczone wtedy
i tylko wtedy, gdy $p_C$ jest ograniczona na każdym
zbiorze ograniczonym w $E$.
\end{lemat}
\begin{proof}
Własności seminormy wynikają z absolutnej wypukłości i
pochłaniania. Jeżeli $B\subseteq tC$, to $p_C(x)\le t$
dla $x\in B$. W drugą stronę, jeśli
$M=\sup_{x\in B}p_C(x)<\infty$, wybieramy $t>M$.
Dla każdego $x\in B$ definicja kresu dolnego zapewnia
$x\in s_x C$ dla pewnego $s_x<t$. Ponieważ $C$ jest
zrównoważony, $s_xC\subseteq tC$, stąd $B\subseteq tC$.
\end{proof}

\begin{twierdzenie}[Kryteria bornologiczności]\label{thm:bornological-equiv}
Niech $E$ będzie lokalnie wypukłą TVS Hausdorffa.
Następujące warunki są równoważne:
\begin{enumerate}[label=\textup{(\roman*)}]
\item $E$ jest bornologiczna;
\item każda seminorma $p:E\to[0,\infty)$, która jest
ograniczona na każdym ograniczonym zbiorze $B\subseteq E$,
jest ciągła;
\item dla każdej lokalnie wypukłej TVS Hausdorffa $F$
każdy liniowy operator $T:E\to F$, który odwzorowuje
zbiory ograniczone w $E$ na zbiory ograniczone w $F$,
jest ciągły.
\end{enumerate}
\end{twierdzenie}
\begin{proof}
(i)$\Rightarrow$(ii). Dla seminormy $p$ jak w (ii)
zbiór $C=\{x:p(x)\le1\}$ jest absolutnie wypukły.
Dla każdego ograniczonego $B$ mamy $M_B=\sup_Bp<\infty$,
więc $B\subseteq tC$ przy dowolnym $t>\max\{1,M_B\}$.
Bornologiczność oznacza, że $C$ jest otoczeniem zera,
co daje ciągłość $p$.

(ii)$\Rightarrow$(iii). Niech $T$ będzie operatorem
odwzorowującym zbiory ograniczone na ograniczone.
Każda ciągła seminorma $q$ na $F$ jest ograniczona
na ograniczonych podzbiorach $F$. Zatem $q\circ T$
jest seminormą ograniczoną na wszystkich ograniczonych
podzbiorach $E$. Z (ii) jest ciągła. Ponieważ topologię
$F$ wyznaczają ciągłe seminormy, operator $T$ jest ciągły.

(iii)$\Rightarrow$(i). Niech $C\subseteq E$ będzie
absolutnie wypukły i pochłaniający wszystkie zbiory
ograniczone. Jego funkcjonał Minkowskiego $p_C$
jest seminormą ograniczoną na każdym takim zbiorze.
Rozważmy przestrzeń $F$, której zbiorem punktów jest $E$,
a topologia jest generowana przez seminormy ciągłe
pierwotnej topologii oraz dodatkowo przez $p_C$.
Jest to lokalnie wypukła TVS Hausdorffa. Identyczność
$T:(E,\text{topologia pierwotna})\to F$ jest liniowa
i odwzorowuje zbiory ograniczone na ograniczone:
wszystkie seminormy wyznaczające topologię $F$
pozostają ograniczone na każdym ograniczonym $B\subseteq E$.
Z (iii) operator $T$ jest ciągły, czyli $p_C$ jest
ciągły w pierwotnej topologii. Z definicji funkcjonału
Minkowskiego wynika $\{x:p_C(x)<1\}\subseteq C$,
zatem $C$ jest otoczeniem zera. Dowiedliśmy (i).
\end{proof}

\begin{twierdzenie}[Przestrzenie unormowane są bornologiczne]
\label{thm:normed-bornological}
Każda przestrzeń unormowana nad $\K$, także niezupełna,
jest bornologiczna.
\end{twierdzenie}
\begin{proof}
Niech $E$ będzie przestrzenią unormowaną, a $C\subseteq E$
zbiorem absolutnie wypukłym, który pochłania wszystkie
zbiory ograniczone. W szczególności pochłania normową
kulę jednostkową $B=\{x:\|x\|\le1\}$, to znaczy
$B\subseteq tC$ dla pewnego $t>0$. Wtedy
$\{x:\|x\|<1/t\}\subseteq C$, więc $C$ jest
otoczeniem zera.
\end{proof}

\begin{przyklad}[Bornologiczna, ale niebeczkowa]
Przestrzeń $c_{00}$ z normą supremum jest bornologiczna
na mocy Twierdzenia~\ref{thm:normed-bornological}, lecz
nie jest beczkowa, jak pokazaliśmy w
Przykładzie~\ref{ex:c00-not-barrelled}.
Zatem bornologiczność nie wystarcza do zastosowania
twierdzenia Banacha--Steinhausa.
\end{przyklad}

\begin{uwaga}[Dwie różne zasady ciągłości]
Beczkowość pozwala przejść od \emph{punktowej ograniczoności
rodziny ciągłych operatorów} do jej równociągłości.
Bornologiczność pozwala zaś przejść od \emph{ograniczoności
pojedynczego operatora liniowego} (przenoszenia zbiorów
ograniczonych na ograniczone) do jego ciągłości.
Te dwa pojęcia mają różne cele i nie są zamienne.
\end{uwaga}

\subsection{Powrót do topologii Mackeya i przestrzeni dystrybucji}

\begin{twierdzenie}[Przestrzenie beczkowe są przestrzeniami Mackeya]
\label{thm:barrelled-mackey}
Niech $E$ będzie beczkową lokalnie wypukłą TVS Hausdorffa
z topologią $\rho$, a $E'$ jej przestrzenią dualną topologiczną.
Wówczas $\rho=\tau(E,E')$, to znaczy pierwotna topologia
jest topologią Mackeya odpowiadającą parze dualnej $(E,E')$.
\end{twierdzenie}
\begin{proof}
Topologia $\rho$ jest zgodna z $(E,E')$, więc z Twierdzenia~\ref{thm:mackey-arens} mamy
$\rho\subseteq\tau(E,E')$.
Niech $K\subseteq E'$ będzie absolutnie wypukłym zbiorem
zwartym w $\sigma(E',E)$. Dla każdego $x\in E$
zbiór $\{f(x):f\in K\}$ jest zwarty w $\K$, a zatem
ograniczony. Twierdzenie~\ref{thm:barrelled-char} daje
równociągłość $K$. Istnieje więc otoczenie $U$ zera,
na którym $p_K(x)=\sup_{f\in K}|f(x)|\le1$.
To dowodzi ciągłości $p_K$ względem $\rho$.
Wszystkie seminormy generujące $\tau(E,E')$ są więc
$\rho$-ciągłe, skąd $\tau(E,E')\subseteq\rho$.
\end{proof}

\begin{wniosek}[Topologia Mackeya przestrzeni Banacha]\label{cor:banach-mackey}
Niech $E$ będzie przestrzenią Banacha, a $E'$ jej
przestrzenią dualną topologiczną. Wówczas
$\tau(E,E')$ jest topologią normową na $E$.
\end{wniosek}
\begin{proof}
Przestrzeń Banacha jest beczkowa, zatem stosujemy
Twierdzenie~\ref{thm:barrelled-mackey}.
\end{proof}

\begin{przyklad}[Aproksymacja dystrybucji Diraca]\label{ex:delta-approx-barrel}
Niech $E=C^\infty([0,1])$ z seminormami
$p_m(\varphi)=\max_{0\le j\le m}\sup_{x\in[0,1]}|\varphi^{(j)}(x)|$,
$m\in\N_0$. Jest to przestrzeń Frécheta, a więc beczkowa.
Dla $n\ge1$ określmy
$T_n(\varphi)=n\int_{[0,1/n]}\varphi(x)\,\mathrm d\lambda(x)$.
Każdy $T_n$ jest ciągłym funkcjonałem liniowym i
$|T_n(\varphi)|\le p_0(\varphi)$, zatem rodzina
$\{T_n:n\ge1\}$ jest równociągła. Ponieważ
$|T_n(\varphi)-\varphi(0)|\le
\sup_{x\in[0,1/n]}|\varphi(x)-\varphi(0)|\to0$,
mamy $T_n(\varphi)\to\varphi(0)$. W późniejszej teorii
dystrybucji funkcjonał $\delta_0(\varphi)=\varphi(0)$
będzie przykładem dystrybucji Diraca; granica zachodzi
punktowo na $E$, a dzięki
Twierdzeniu~\ref{thm:equicont-compact-conv} również
jednostajnie na każdym zwartym zbiorze funkcji testowych
w topologii $E$.
\end{przyklad}

\begin{uwaga}[Granica zakresu obecnego rozdziału]
Nie należy utożsamiać zwartości i ograniczoności w
$C^\infty([0,1])$ ani przenosić automatycznie jej własności
na przestrzeń $\mathcal D(\Omega)$ z topologią LF.
Precyzyjne omówienie zbiorów ograniczonych w przestrzeniach
funkcji testowych nastąpi po konstrukcji granic induktywnych.
\end{uwaga}
